\section{Annihilator, Torsion and Projectivity of Free Modules}

\begin{definition}
    Let $M$ be an $R$-module, and let $m \in M$. The cyclic submodule generated by $m$ is
    \[
        Rm = \langle m \rangle = \{\, rm \mid r \in R \,\}.
    \]
    The \emph{annihilator} of $m$ in $R$ is
    \[
        \operatorname{Ann}_R(m) = \{\, r \in R \mid rm = 0 \,\}.
    \]
    It is a left ideal of $R$.
\end{definition}

\begin{proposition}
    For every $m \in M$, there is an isomorphism of left $R$-modules
    \[
        R/\operatorname{Ann}_R(m) \cong Rm.
    \]
    In particular, every cyclic module is isomorphic to a quotient of $R$.
\end{proposition}

\begin{definition}
    An element $m \in M$ is called
    \begin{itemize}
        \item \textbf{torsion-free} if $\operatorname{Ann}_R(m) = 0$;
        \item \textbf{torsion} if $\operatorname{Ann}_R(m) \neq 0$.
    \end{itemize}
    A module is called \emph{torsion-free} if every nonzero element is torsion-free.
\end{definition}

\begin{definition}
    The \emph{torsion submodule} of $M$ is
    \[
        \operatorname{Tor}(M) = \{\, m \in M \mid \operatorname{Ann}_R(m) \neq 0 \,\}.
    \]
\end{definition}

\begin{lemma}
    If $R$ is an integral domain, then the torsion part of $M$ is a submodule of $M$.
\end{lemma}

\begin{proof}
    We verify the submodule criteria.

    \emph{Non-empty.}
    Since $r\cdot 0_M=0_M$ for every $r\in R$ and $R\neq 0$,
    we have $\operatorname{Ann}_R(0_M)=R\neq 0$, so $0_M\in\operatorname{Tor}(M)$.

    \emph{Closed under addition.}
    Let $m,n\in\operatorname{Tor}(M)$.
    Then there exist $\lambda,\mu\in R\setminus\{0\}$ such that $\lambda m=0$ and $\mu n=0$.
    Hence
    \[
        (\lambda\mu)(m+n)=\mu(\lambda m)+\lambda(\mu n)=0.
    \]
    Since $R$ is an integral domain and $\lambda,\mu\neq 0$, we have $\lambda\mu\neq 0$,
    so $\operatorname{Ann}_R(m+n)\neq 0$ and $m+n\in\operatorname{Tor}(M)$.

    \emph{Closed under scalar multiplication.}
    Let $m\in\operatorname{Tor}(M)$ and $r\in R$.
    Choose $\lambda\in R\setminus\{0\}$ with $\lambda m=0$. Then
    \[
        \lambda(rm)=r(\lambda m)=0,
    \]
    so $\operatorname{Ann}_R(rm)\neq 0$ and $rm\in\operatorname{Tor}(M)$.

    Therefore $\operatorname{Tor}(M)$ is a submodule of $M$.
\end{proof}

\begin{example}
    (1) The element $0 \in M$ is torsion.

    (2) Let $R$ be a commutative ring with $1$. Then $R$ is an integral domain if and only if the left regular module ${}_R R$ is torsion-free. In this case $\operatorname{Tor}({}_R R)=0$ and $\operatorname{Tor}(R^n)=0$ for all $n \ge 1$.

    (3) If $I$ is a nonzero ideal of $R$, then $\operatorname{Tor}(R/I)=R/I$ when $R/I$ is regarded as an $R$-module. If we regard $R/I$ as an $R/I$-module, then it is torsion-free if and only if $I$ is a prime ideal.

    (4) More generally, let $I$ be any nonzero ideal of $R$, and let $M$ be an $R$-module. Then $\operatorname{Tor}(M/IM)=M/IM$ when $M/IM$ is regarded as an $R$-module.

    (5) Even when $I$ is prime, $M/IM$ need not be torsion-free as an $R/I$-module. For $\overline{m}\in M/IM$, one has $\operatorname{Ann}_{R/I}(\overline{m})\neq 0$ if and only if there exists $r\in R\setminus I$ with $rm\in IM$. Thus primeness of $I$ does not rule out nonzero operators in $R/I$ annihilating nonzero classes in $M/IM$.

    For instance, let $R=k[x,y]$, let $I=(y)$, and let $M=R/(x)\cong k[y]$. Then $IM=(y)$ and $M/IM\cong k$, while $x\cdot M=0$ but $x\notin I$; hence $x+I\in R/I$ is nonzero and annihilates all of $M/IM$.
    If instead $I$ is maximal, then $R/I$ is a field, nonzero scalars act invertibly, and $M/IM$ is always torsion-free over $R/I$.
    Item~\textup{(3)} is the case $M=R$.
\end{example}

\begin{remark}
    Let $R$ be a commutative ring with $1$, let $M$ be an $R$-module, and fix $r\in R$. Then the map
    \[
        f\colon M \to M, \qquad f(m)=rm,
    \]
    is an $R$-module homomorphism. If moreover $M$ is torsion-free and $r\neq 0$, then $f$ is a monomorphism.
\end{remark}

\begin{lemma}
    If $R$ is an integral domain and $M$ is an $R$-module, then
    \begin{enumerate}
        \item $\operatorname{Tor}(\operatorname{Tor}(M)) = \operatorname{Tor}(M)$.
        \item $\operatorname{Tor}(M/\operatorname{Tor}(M)) = 0$.
    \end{enumerate}
\end{lemma}

\begin{proof}\leavevmode
    \begin{enumerate}
        \item
              Since $\operatorname{Tor}(M)$ is a submodule of $M$, its torsion submodule is
              \[
                  \operatorname{Tor}(\operatorname{Tor}(M))
                  = \{\, m\in\operatorname{Tor}(M) \mid \operatorname{Ann}_R(m)\neq 0 \,\}.
              \]
              Every $m\in\operatorname{Tor}(M)$ already satisfies $\operatorname{Ann}_R(m)\neq 0$,
              so $\operatorname{Tor}(\operatorname{Tor}(M)) \subseteq \operatorname{Tor}(M)$.
              Conversely, $\operatorname{Tor}(M)\subseteq\operatorname{Tor}(\operatorname{Tor}(M))$ because the annihilator of $m$ in $R$ does not depend on the ambient module:
              for any $m\in\operatorname{Tor}(M)$, its annihilator computed in the submodule $\operatorname{Tor}(M)$ is the same as $\operatorname{Ann}_R(m)\neq 0$,
              so $m\in\operatorname{Tor}(\operatorname{Tor}(M))$.
              Hence $\operatorname{Tor}(\operatorname{Tor}(M))=\operatorname{Tor}(M)$.

        \item
              Let $\pi\colon M\to M/\operatorname{Tor}(M)$ be the canonical projection.
              Write $\overline{m}=\pi(m)=m+\operatorname{Tor}(M)$.
              Suppose $\overline{m}\in\operatorname{Tor}(M/\operatorname{Tor}(M))$,
              i.e.\ there exists $\lambda\in R\setminus\{0\}$ such that
              \[
                  \lambda\,\overline{m} = \overline{0} \quad\text{in } M/\operatorname{Tor}(M).
              \]
              This means $\lambda m\in\operatorname{Tor}(M)$,
              so there exists $\mu\in R\setminus\{0\}$ with $\mu(\lambda m)=0$, i.e.
              \[
                  (\mu\lambda)\,m = 0.
              \]
              Since $R$ is an integral domain and $\mu,\lambda\neq 0$, we have $\mu\lambda\neq 0$,
              hence $m\in\operatorname{Tor}(M)$, so $\overline{m}=\overline{0}$.
              Therefore $\operatorname{Tor}(M/\operatorname{Tor}(M))=0$.
    \end{enumerate}
\end{proof}

\begin{definition}
    Let $J$ be a subscript set, denote $R^J=\prod_{j\in J} R$.$R^{(J)}=\bigoplus_{j\in J} R$.

    Let $M$ be an $R$-module. Let $x_j \in M$ where $j \in J$ be a family of elements of $M$. Then the submodule generated by $(x_j)_{j\in J}$ is denoted by $\sum_{j\in J} Rx_j$.
\end{definition}

\begin{corollary}
    We have a morphism of $R$-modules $\phi: R^{(J)} \to M$ such that $\phi(r_j)_{j \in J}=\sum_{j \in J}r_jx_j$ for all $(r_j)_{j \in J} \in R^{(J)}$.
    \begin{itemize}
        \item $\phi$ is injective if and only if $\{x_j\}_{j\in J}$ are linearly independent.
        \item $\phi$ is surjective if and only if $\{x_j\}_{j\in J}$ generate $M$.
        \item $\phi$ is an isomorphism if and only if $M$ is free with basis $\{x_j\}_{j\in J}$.
    \end{itemize}
\end{corollary}

\begin{lemma}
    Let $L$ be a free $R$-module with basis $\{x_j\}_{j\in J}$. For any $R$-module $M$, the map
    \[
        \Phi:\operatorname{Hom}_R(L,M)\to M^J,\qquad
        f\mapsto (f(x_j))_{j\in J}
    \]
    is a bijection.
\end{lemma}

\begin{proof}
    We prove that $\Phi$ is injective and surjective.

    \medskip
    \noindent\textbf{Injectivity.}
    Let $f,g\in \operatorname{Hom}_R(L,M)$ and suppose that
    \[
        \Phi(f)=\Phi(g).
    \]
    Then for every $j\in J$ we have
    \[
        f(x_j)=g(x_j).
    \]
    Since $\{x_j\}_{j\in J}$ is a basis of $L$, every element $x\in L$ can be written uniquely as a finite sum
    \[
        x=\sum_{j\in J} a_j x_j
    \]
    with $a_j\in R$ and all but finitely many $a_j$ equal to $0$. Hence
    \[
        f(x)=f\!\left(\sum_{j\in J} a_j x_j\right)
        =\sum_{j\in J} a_j f(x_j)
        =\sum_{j\in J} a_j g(x_j)
        =g\!\left(\sum_{j\in J} a_j x_j\right)
        =g(x).
    \]
    Thus $f=g$, so $\Phi$ is injective.

    \medskip
    \noindent\textbf{Surjectivity.}
    Let $(m_j)_{j\in J}\in M^J$ be given. We define a map
    \[
        f:L\to M
    \]
    by declaring that for
    \[
        x=\sum_{j\in J} a_j x_j \qquad (a_j\in R,\ \text{all but finitely many } a_j=0),
    \]
    we set
    \[
        f(x):=\sum_{j\in J} a_j m_j.
    \]
    This is well defined because each element of $L$ has a unique expression as a finite linear combination of the basis elements.

    Now let
    \[
        x=\sum_{j\in J} a_j x_j,\qquad
        y=\sum_{j\in J} b_j x_j,\qquad
        r\in R.
    \]
    Then
    \[
        f(x+y)
        =f\!\left(\sum_{j\in J} (a_j+b_j)x_j\right)
        =\sum_{j\in J}(a_j+b_j)m_j
        =\sum_{j\in J}a_jm_j+\sum_{j\in J}b_jm_j
        =f(x)+f(y),
    \]
    and
    \[
        f(rx)
        =f\!\left(\sum_{j\in J} (ra_j)x_j\right)
        =\sum_{j\in J} (ra_j)m_j
        =r\sum_{j\in J} a_jm_j
        =rf(x).
    \]
    So $f$ is an $R$-module homomorphism. Moreover, for every $j\in J$,
    \[
        f(x_j)=m_j.
    \]
    Therefore
    \[
        \Phi(f)=(m_j)_{j\in J}.
    \]
    Thus $\Phi$ is surjective.

    Since $\Phi$ is both injective and surjective, it is a bijection.
\end{proof}

\begin{proposition}
    Let $L$ be a free $R$-module.
    \begin{enumerate}
        \item Let $\pi: X\to Y$ be an epimorphism of $R$-modules. For any $R$-module homomorphism $\varphi: L \to Y$ there exists a morphism $\tilde{\varphi}: L \to X$ such that the following diagram commutes:
              \begin{center}
                  \begin{tikzcd}
                      & L \arrow[d, "\varphi"] \arrow[ld, "\tilde{\varphi}"'] &   \\
                      X \arrow[r, "\pi"'] & Y \arrow[r]                                           & 0
                  \end{tikzcd}
              \end{center}
        \item In particular, for any $R$-module $M$, any epimorphism $\pi: M \to L$ has a section, that is, there exists a morphism $\sigma: L \to M$ such that $\pi \circ \sigma = \operatorname{id}_L$. In that case $\sigma$ is injective, $\operatorname{Im}(\sigma)\cong L$, and
              \[
                  M = \ker \pi \oplus \operatorname{Im}(\sigma).
              \]
              \begin{center}
                  \begin{tikzcd}
                      0 \arrow[r] & \ker \pi \arrow[r] & M \arrow[r, two heads, shift left, "\pi"] & L \arrow[l, dashed, shift left, "\sigma"] \arrow[r] & 0
                  \end{tikzcd}
              \end{center}
    \end{enumerate}
\end{proposition}

\begin{proof}
    \begin{enumerate}
        \item Since $L$ is free, let $\{x_j\}_{j\in J}$ be a basis of $L$.

              For each $j\in J$, since $\pi:X\to Y$ is surjective, there exists an element
              \[
                  \tilde{x}_j \in X
              \]
              such that
              \[
                  \pi(\tilde{x}_j)=\varphi(x_j).
              \]
              Because $L$ is free with basis $\{x_j\}_{j\in J}$, there exists a unique $R$-module homomorphism
              \[
                  \tilde{\varphi}:L\to X
              \]
              satisfying
              \[
                  \tilde{\varphi}(x_j)=\tilde{x}_j
                  \qquad \text{for all } j\in J.
              \]
              We now check that $\pi\circ \tilde{\varphi}=\varphi$. For each basis element $x_j$,
              \[
                  (\pi\circ \tilde{\varphi})(x_j)
                  =\pi(\tilde{\varphi}(x_j))
                  =\pi(\tilde{x}_j)
                  =\varphi(x_j).
              \]
              Hence $\pi\circ \tilde{\varphi}$ and $\varphi$ agree on a basis of $L$, so they are equal as $R$-module homomorphisms. Therefore the diagram commutes.

        \item Apply part (1) with
              \[
                  X=M,\qquad Y=L,\qquad \varphi=\operatorname{id}_L.
              \]
              Since $\pi:M\to L$ is an epimorphism, there exists an $R$-module homomorphism
              \[
                  \sigma:L\to M
              \]
              such that
              \[
                  \pi\circ \sigma=\operatorname{id}_L.
              \]
              Thus $\sigma$ is a section of $\pi$.

              We first show that $\sigma$ is injective. If $\sigma(\ell)=0$ for some $\ell\in L$, then
              \[
                  \ell = (\pi\circ \sigma)(\ell)=\pi(0)=0.
              \]
              Hence $\ker \sigma=0$, so $\sigma$ is injective.

              Next, since $\sigma$ is injective, it induces an isomorphism
              \[
                  L \xrightarrow{\sim} \operatorname{Im}(\sigma),
                  \qquad \ell \mapsto \sigma(\ell).
              \]
              Thus $\operatorname{Im}(\sigma)\cong L$.

              It remains to prove that
              \[
                  M=\ker \pi \oplus \operatorname{Im}(\sigma).
              \]

              First, let $m\in M$. Then
              \[
                  m=\bigl(m-\sigma(\pi(m))\bigr)+\sigma(\pi(m)).
              \]
              Now
              \[
                  \pi\bigl(m-\sigma(\pi(m))\bigr)
                  =\pi(m)-(\pi\circ \sigma)(\pi(m))
                  =\pi(m)-\pi(m)=0,
              \]
              so
              \[
                  m-\sigma(\pi(m))\in \ker \pi,
              \]
              while clearly
              \[
                  \sigma(\pi(m))\in \operatorname{Im}(\sigma).
              \]
              Hence
              \[
                  M=\ker \pi+\operatorname{Im}(\sigma).
              \]

              Finally, suppose
              \[
                  z\in \ker \pi \cap \operatorname{Im}(\sigma).
              \]
              Then $z=\sigma(\ell)$ for some $\ell\in L$, and also $\pi(z)=0$. Therefore
              \[
                  0=\pi(z)=\pi(\sigma(\ell))=(\pi\circ \sigma)(\ell)=\ell.
              \]
              So $z=\sigma(\ell)=0$. Thus
              \[
                  \ker \pi \cap \operatorname{Im}(\sigma)=0.
              \]
              Therefore
              \[
                  M=\ker \pi \oplus \operatorname{Im}(\sigma).
              \]
    \end{enumerate}
\end{proof}

\begin{exercise}
    Let $L$ be a free $R$-module. Let $P$ be a direct summand of $L$. Show that whenever
    \[
        \pi:X\to Y
    \]
    is an epimorphism and
    \[
        \varphi:P\to Y
    \]
    is a morphism, there exists a morphism
    \[
        \tilde{\varphi}:P\to X
    \]
    such that
    \[
        \varphi=\pi\circ \tilde{\varphi}.
    \]
\end{exercise}

\begin{proof}
    Since $P$ is a direct summand of $L$, there exist $R$-module homomorphisms
    \[
        i:P\to L
        \qquad\text{and}\qquad
        p:L\to P
    \]
    such that
    \[
        p\circ i=\operatorname{id}_P.
    \]
    (Here $i$ is the inclusion and $p$ is the projection.)

    Consider the composite
    \[
        \psi:=\varphi\circ p:L\to Y.
    \]
    Since $L$ is free and $\pi:X\to Y$ is an epimorphism, by the lifting property of free modules there exists an $R$-module homomorphism
    \[
        \widetilde{\psi}:L\to X
    \]
    such that
    \[
        \pi\circ \widetilde{\psi}=\psi.
    \]

    Now define
    \[
        \widetilde{\varphi}:=\widetilde{\psi}\circ i:P\to X.
    \]
    Then
    \[
        \pi\circ \widetilde{\varphi}
        =\pi\circ \widetilde{\psi}\circ i
        =\psi\circ i
        =\varphi\circ p\circ i
        =\varphi\circ \operatorname{id}_P
        =\varphi.
    \]
    Therefore there exists a morphism $\widetilde{\varphi}:P\to X$ such that
    \[
        \varphi=\pi\circ \widetilde{\varphi}.
    \]
\end{proof}

